% =====================================================================
%  3 -- STATISTICAL EQUILIBRIUM: NSE AND QSE  (source node S9).
%
%  Source: tier2.md Part III -- header (3019-3030), III.1 (3034-3078),
%          III.2 (3082-3134), III.3 (3138-3186), III.4 (3190-3214),
%          III.5 (3218-3242), III.6 (3246-3329), III.7 (3333-3351),
%          III.8 (3355-3406), III.9 (3410-3452), III.11 (3520-3551),
%          III.12 result (1) (3555-3563), III.13 (3613-3635),
%          III.14 self-check (3639-3656).
%
%  Excluded and re-homed FORWARD into section 4, because both are
%  defined against kappa and the physics arc puts equilibrium first:
%          III.10 the two departure diagnostics (3456-3516) ->
%                 sec:departure;
%          III.12 results (2) and (3) -- the absent plateau and the
%                 empty maskable set (3565-3609) -> sec:negresults.
%          III.14 self-check items 7 and 8 travel with them.
%  The III.13 oracle table is kept whole here (it is the qse/ oracle
%  list) with its two diagnostics rows cross-referenced forward.
% =====================================================================

\section{Statistical equilibrium: NSE and QSE}
\label{part:equilibrium}

\begin{repopointer}
\textbf{Files:} \code{qse/coeffs.py}, \code{qse/solver.py},
\code{qse/diagnostics.py}, \code{configs/qse\_groups.yaml}. \quad
\textbf{Oracles:} \code{tests/test\_qse\_coeffs.py},
\code{tests/test\_qse\_solver.py} --- \S\ref{sec:codes9}. \quad
\textbf{Notebook:} \code{06-qse-bridges.ipynb}.
\end{repopointer}

\S\ref{part:setting} said equilibrium is what makes silicon burning tractable.
This section makes it \emph{computable}: an independent Saha solver, a cluster
generalisation, and --- in \S\ref{sec:departure}, once \kap{} is available ---
two departure diagnostics. Without it, ``this reaction is equilibrated'' has no
measurable meaning, and the kill-test has no reference.

\begin{aside}
\textbf{Why this comes before the flux algebra.} In the source document this
node sits \emph{after} fluxes and cancellation. The physics runs the other way:
the hinge identity of the whole tier is $\kap = |\tanh(\aff/2kT)|$
(\S\ref{sec:kappa-affinity}), and the affinity \aff{} is built from the chemical
potentials that \S\ref{sec:saha}'s Saha coefficients produce. Deriving
equilibrium first turns three forward references into backward ones. The price
is that the two diagnostics defined \emph{against} \kap{} ---
\S\ref{sec:departure} and the mask verdicts of \S\ref{sec:negresults} --- travel
forward into \S\ref{part:fluxes} with it.
\end{aside}

\textbf{The independence matters.} \code{qse/} deliberately uses \textbf{the
same nuclear inputs} as pynucastro (Rauscher partition functions,
\code{Nucleus.nucbind}, \code{Nucleus.A\_nuc}, \code{Nucleus.spin\_states}) but
an entirely separate solver. Shared data $+$ separate algorithm $=$ the
cross-check tests the \emph{solver}, not the data. If it used its own masses
too, agreement would prove nothing about either.

% =====================================================================
\subsection{Equilibrium from the ensemble}
\label{sec:ensemble}

\subsubsection{The distribution}
\label{sec:distribution}

For a species $i$ treated as an ideal, non-degenerate, non-relativistic gas with
ground-state spin degeneracy $(2J_i + 1)$ and internal partition function
$G_i(T)$, the grand-canonical number density is (the standard NSE starting
point: \citealp{clifford1965}; \citealp{hwe1985}; chemical-potential form as in
\citealp{seitenzahl2009})
%
\begin{equation}
n_i \;=\; (2J_i+1)\,G_i(T)\left(\frac{m_i kT}{2\pi\hbar^2}\right)^{3/2}
\exp\!\left(\frac{\mu_i - m_i c^2}{kT}\right),
\label{eq:ndensity}
\end{equation}
%
with $\mu_i$ the full chemical potential (rest mass included). Define the
\textbf{kinetic} chemical potential $u_i \equiv \mu_i - m_i c^2$; this is what
the code's \code{u\_p}, \code{u\_n} are, in MeV.

\paragraph{Is Maxwell--Boltzmann legitimate here?} It must be checked, because
the \emph{electrons} in this plasma are strongly degenerate (Tier~0 \S0.1). The
criterion is $n_i \lambda_i^3 \ll 1$ with $\lambda = h/\sqrt{2\pi m kT}$. At box
centre ($\Tn = 4$, $\rho = 10^{8}$\,\gcc) for $A = 28$:
$m = 4.65\times10^{-23}$\,g, $kT = 5.52\times10^{-7}$\,erg, so
$\lambda = 5.2\times10^{-13}$\,cm, $\lambda^3 = 1.4\times10^{-37}$\,cm$^3$, and
at $X = 1$, $n = \rho/(A m_u) = 2.2\times10^{30}$\,cm$^{-3}$. Hence
%
\begin{equation}
n\lambda^3 \;\approx\; 3\times10^{-7} \;\lll\; 1.
\label{eq:nlambda3}
\end{equation}
\derivedhere

Six and a half orders of margin. Nuclei are classical here; electrons are not,
because $m_e$ is $5\times10^{4}$ times smaller and $\lambda_e^3$ correspondingly
$10^{7}$ times larger. \textbf{The same plasma is simultaneously classical in
its nuclei and degenerate in its electrons}, which is why the electron sector
needs a Fermi treatment (Tier~1 \S V.4) while the Saha algebra can stay
Boltzmann.

\subsubsection{The equilibrium condition}
\label{sec:eqcondition}

Chemical equilibrium for reaction $j$ is $\sum_i \nu_{ij}\mu_i = 0$ (in the full
$\mu$). Since the reaction graph is connected through the light particles
(\S\ref{sec:graphK}) and every nucleus is reachable from free nucleons by strong
reactions, all such conditions collapse to a single family:
%
\begin{equation}
\mu_i \;=\; Z_i\,\mu_p + N_i\,\mu_n .
\label{eq:muchain}
\end{equation}

% =====================================================================
\subsection{Where the binding energy comes from}
\label{sec:binding}

Substituting $\mu = u + mc^2$:
%
\begin{equation}
u_i + m_ic^2 \;=\; Z_i(\up + m_pc^2) + N_i(\un + m_nc^2)
\label{eq:musub}
\end{equation}
%
\begin{equation}
\Longrightarrow\quad u_i \;=\; Z_i \up + N_i \un +
  \underbrace{\bigl[Z_i m_p + N_i m_n - m_i\bigr]c^2}_{\textstyle B_i}.
\label{eq:bindingappears}
\end{equation}

\begin{invariant}
\textbf{The binding energy appears as the rest-mass bookkeeping of the
equilibrium condition, not as an extra assumption.} This is worth doing once by
hand, because it is the reason \code{nucbind} shows up in a
\emph{chemical-potential} formula and the reason the equilibrium distribution
favours the most-bound nucleus: $B_i/kT$ is an exponential reward for binding.
\end{invariant}

At $\Tn = 4$, $kT = 0.345$\,MeV, and $B(\nuc{Ni}{56}) = 484$\,MeV, so
$B/kT \approx 1400$. The Saha exponent is enormous; what keeps abundances finite
is the competition against the $-\ln\rho$ and $(Z\up + N\un)/kT$ terms, with
\up{} and \un{} both strongly negative. Measured, at $\rho = 10^{8}$\,\gcc:
\derivedhere

\begin{center}
\begin{tabularx}{\textwidth}{@{}L{2.4cm} L{2.4cm} L{2.4cm} Y@{}}
\toprule
\textbf{$(\Tn, \Ye)$} & \textbf{\up{} [MeV]} & \textbf{\un{} [MeV]} &
\textbf{dominant species} \\
\midrule
$(4, 0.50)$   & $-5.09$ & $-12.41$ & \nuc{Ni}{56} ($X = 0.94$) \\
$(4, 0.47)$   & $-7.47$ & $-10.14$ & \nuc{Fe}{56}, \nuc{Fe}{54}, \nuc{Ni}{58} \\
$(3.5, 0.47)$ & $-7.43$ & $-10.15$ & \nuc{Fe}{56}, \nuc{Fe}{54}, \nuc{Ni}{58} \\
$(5, 0.47)$   & $-7.58$ & $-10.10$ & \nuc{Fe}{56}, \nuc{Fe}{54}, \nuc{Ni}{58} \\
\bottomrule
\end{tabularx}
\end{center}

Two things to read off. The $\Ye = 0.47$ rows are \S\ref{sec:fepeak}'s endpoint
claim, measured: drop \Ye{} three points below 0.5 and the equilibrium abandons
\nuc{Ni}{56} for the neutron-richer Fe-peak species. And the \emph{difference}
$\up - \un$ is the free-nucleon ratio, exactly --- since $B_p = B_n = 0$ and the
mass/spin factors are $\approx 1$,
%
\begin{equation}
\frac{X_p}{X_n} \;\approx\; \exp\!\frac{\up - \un}{kT},
\label{eq:xpxn}
\end{equation}
%
which the solver confirms: at $\Ye = 0.50$, $\up - \un = 7.32$\,MeV $= 21\,kT$
and $X_p/X_n = 1.7\times10^{9}$ ($X_p = 6.3\times10^{-4}$,
$X_n = 3.8\times10^{-13}$); at $\Ye = 0.47$, $\up - \un = 2.67$\,MeV
$= 7.7\,kT$ and the ratio is $2.3\times10^{3}$. Note the free nucleons are
\emph{not} balanced even at $\Ye = 0.5$ --- with \nuc{Ni}{56} holding 94\% of
the mass there is almost no neutron the composition needs to leave free, so
$\up - \un$ runs to tens of $kT$. The lesson for reading a solver trace is that
$\up - \un$ is a steep, directly interpretable function of \Ye, not that it is
small.\derivedhere

\begin{warn}
\textbf{Do not read the solver's \code{init=(-3.5, -15.0)} as typical values.}
It is a seed chosen to sit safely inside the convergence basin, not a converged
answer, and nothing in the table above lands near it --- \up{} is 1.5--4\,MeV
too high and \un{} 2--5\,MeV too low. The arrow in the original claim runs
backwards: the seed is ``somewhere Newton reliably converges \emph{from}'',
which is a numerical property, and inferring physics from it is exactly the
mistake.
\end{warn}

% =====================================================================
\subsection{The Saha coefficient, term by term}
\label{sec:saha}

Convert to mass fraction: $X_i = n_i m_i/\rho$ with
$m_i = A_{\mathrm{nuc},i}\cdot m_u$.
%
\begin{equation}
\ln X_i = \ln(A_{\mathrm{nuc},i} m_u) - \ln\rho + \ln(2J_i+1) + \ln G_i
+ \tfrac32\ln\!\frac{m_i kT}{2\pi\hbar^2} + \frac{Z_i\up + N_i\un + B_i}{kT}.
\label{eq:lnxi}
\end{equation}
%
Collect the two mass terms:
$\ln(A_{\rm nuc} m_u) + \tfrac32\ln(A_{\rm nuc} m_u) =
\mathbf{2.5}\cdot\ln(A_{\rm nuc} m_u)$. Hence
%
\begin{equation}
\boxed{\ \ln X_i \;=\;
\underbrace{\tfrac52\ln(A_{\mathrm{nuc},i}m_u) + \ln(2J_i+1) - \ln\rho
+ \tfrac32\ln\!\frac{kT}{2\pi\hbar^2} + \ln G_i(T) + \frac{B_i}{kT}}
_{\textstyle \log C_i(T,\rho)}
\; + \; \frac{Z_i\up + N_i\un}{kT}\ }
\label{eq:sahacoeff}
\end{equation}
%
which is \code{qse/coeffs.py::nse\_log\_coeffs\_batch}, line for line:

\begin{srclisting}
species = 2.5*np.log(inputs.A_nuc * constants.m_u_C18) + np.log(inputs.spin_states)
bind    = inputs.nucbind * inputs.A                     # B/A x A = B
state   = -np.log(rho) + 1.5*np.log(constants.k*T / (2.0*np.pi*constants.hbar**2))
return species[None,:] + state[:,None] + log_pf + bind[None,:]/kT_MeV[:,None]
\end{srclisting}

\paragraph{Three details that are easy to get subtly wrong.}

\begin{enumerate}
\item \textbf{$A_{\rm nuc}$ vs $A$.} The mass prefactor uses \code{A\_nuc} (the
  \emph{actual} mass in amu: 27.9769 for \nuc{Si}{28}), the binding term uses
  the integer \code{A} (nucleon count). Using $A$ for both would introduce a
  factor $\exp(2.5\cdot\ln(A/A_{\rm nuc})) \approx 1.002$ --- a 0.2\% error,
  invisible in a plot and fatal to a $10^{-12}$ cross-check.
\item \textbf{$G$ vs $(2J{+}1)G$.} Tier~1 \S III.3 flags this normalisation trap
  on the rate side; the same trap exists here, and the code splits the two
  explicitly (\code{spin\_states} and \code{log\_pf} as separate terms) rather
  than relying on a library convention.
\item \textbf{\code{EXP\_CLIP = 500}.} The Saha exponent is clipped before
  exponentiation, matching pynucastro (\code{nse\_network.py:212}).
  $\exp(500) \approx 10^{217}$, comfortably inside float64's
  $1.8\times10^{308}$ --- so the clip prevents overflow during Newton excursions
  from a bad guess \emph{without} affecting any converged answer, where
  $X \le 1$.
\end{enumerate}

\paragraph{The cross-check.} \code{tests/test\_qse\_coeffs.py} compares these
mass fractions at \emph{fixed} $(\up, \un)$ against pynucastro's
\code{\_nucleon\_fraction\_nse}, $\le 10^{-12}$ rel. Fixing $u$ removes the
solver from the comparison entirely: this test isolates the coefficient
construction. Only then does \code{test\_qse\_solver.py} compare converged
solutions. \textbf{Layered oracles again} --- the same discipline as
\code{test\_conservation.py}'s two layers (\S\ref{sec:oracles7}).

% =====================================================================
\subsection{The two constraints}
\label{sec:twoconstraints}

$\log C$ determines the \emph{shape}; two constraints fix $(\up, \un)$:
%
\begin{equation}
F_0(\up,\un) = \sum_i X_i - 1 = 0, \qquad
F_1(\up,\un) = \sum_i \frac{Z_i}{A_i}X_i - \Ye = 0 .
\label{eq:nseconstraints}
\end{equation}
%
The second is $\Ye = \sum_i Z_i Y_i = \sum_i (Z_i/A_i) X_i$ (Tier~0 \S II.5).
Note what is \emph{not} a constraint: nothing about neutrons separately, because
baryon number is already carried by $\sum X = 1$. Two constraints, two unknowns
--- the count of \S\ref{sec:regimes}.

\paragraph{And note the convention mix, because \S\ref{sec:saha} just made a
fuss about exactly this distinction.} $X_i$ is built from the \emph{actual} mass
$A_{\rm nuc}$, but $F_1$ uses the \textbf{integer} $A$ (\code{solver.py}:
\code{((Z / A) * X).sum() - ye}). Strictly, if $\sum A_{{\rm nuc},i} Y_i = 1$
then $\Ye = \sum_i Z_i X_i / A_{{\rm nuc},i}$, and using $Z/A$ instead shifts it
by ${\sim}0.1\%$. This is deliberate and correct, for a reason worth stating:
\textbf{\Ye{} here is not a derived quantity being computed accurately, it is a
boundary condition being imposed}, and the convention that must match is the one
used by everything that consumes it --- pynucastro's NSE solver, MESA's
\code{ye}, and the \Ye{} column of the training data all define
$\Ye = \sum(Z_i/A_i)X_i$ with integer $A$. Using $A_{\rm nuc}$ here would make
the solver internally prettier and externally inconsistent, which is the wrong
trade. \S\ref{sec:saha}'s detail (1) is about a quantity being \emph{computed};
this is a quantity being \emph{matched}.

% =====================================================================
\subsection{Newton with an analytic Jacobian}
\label{sec:newton}

Since $\partial X_i/\partial \up = X_i Z_i/kT$ and
$\partial X_i/\partial \un = X_i N_i/kT$, the Jacobian is available in closed
form and costs nothing beyond the residual evaluation:
%
\begin{equation}
J \;=\; \frac{1}{kT}
\begin{pmatrix}
\sum X_iZ_i & \sum X_iN_i\\[2pt]
\sum X_iZ_i^2/A_i & \sum X_iZ_iN_i/A_i
\end{pmatrix}
\label{eq:nsejacobian}
\end{equation}

\begin{srclisting}
J = np.array([[(X*Z).sum(),       (X*N).sum()],
              [(X*Z*Z/A).sum(),   (X*Z*N/A).sum()]]) / kT
\end{srclisting}

\paragraph{Step damping.} \code{\_MAX\_STEP\_MEV = 2.0}: the Newton step is
capped at 2\,MeV in $\infty$-norm. Why a cap is \emph{necessary} rather than
merely prudent falls straight out of \S\ref{sec:monotonicity}'s derivative: when
one species dominates, the compositional variance of $Z/A$ collapses and the
Jacobian becomes nearly singular in the \up{} direction, so an undamped step can
be tens of MeV --- which lands in a region where the exponent clips and the
residual carries no gradient at all. The cap is a trust region; 2\,MeV is
${\sim}6\,kT$ at $\Tn = 4$, so it never obstructs a well-scaled step.

% =====================================================================
\subsection{The monotonicity proof --- why the bisection fallback works}
\label{sec:monotonicity}

\begin{warn}
\warnglyph{}~The fallback is a nested bisection, and it rests on two monotonicity
claims that the code states as comments. Both are provable, and the second is
not obvious.
\end{warn}

\paragraph{Claim 1.} At fixed \up, $\sum X$ is strictly increasing in \un.
%
\begin{equation}
\frac{\partial}{\partial \un}\sum_i X_i \;=\; \frac{1}{kT}\sum_i X_iN_i \;>\;0
\label{eq:claim1}
\end{equation}
%
whenever any species with $N > 0$ has nonzero abundance. Immediate. This
licenses \code{u\_n\_for\_mass}, the inner bisection.

\paragraph{Claim 2.} Along the curve $\sum X = 1$, \Ye{} is strictly increasing
in \up.

Not immediate --- \un{} must move to keep $\sum X = 1$, and it moves
\emph{down}, which pushes \Ye{} up further but also removes mass. Do it
properly. Write $w_i = X_iA_i$, $a_i = Z_i/A_i$, and
%
\begin{equation}
S_Z=\sum X_iZ_i,\quad S_N=\sum X_iN_i,\quad
M_{ZZ}=\sum X_i\frac{Z_i^2}{A_i},\quad
M_{ZN}=\sum X_i\frac{Z_iN_i}{A_i}.
\label{eq:moments}
\end{equation}
%
Implicit differentiation of $\sum X = 1$ gives $\dd\un/\dd\up = -S_Z/S_N$. Then
%
\begin{equation}
kT\,\frac{\dd\Ye}{\dd\up}\bigg|_{\Sigma X=1}
= M_{ZZ} - \frac{S_Z}{S_N}M_{ZN}
= \frac{1}{S_N}\Bigl(M_{ZZ}S_N - M_{ZN}S_Z\Bigr).
\label{eq:dyedup1}
\end{equation}
%
Substitute $Z_i = a_iA_i$, $N_i = (1-a_i)A_i$ so that $M_{ZZ} = \sum w a^2$,
$S_N = \sum w(1-a)$, $M_{ZN} = \sum w a(1-a)$, $S_Z = \sum w a$:
%
\begin{equation}
M_{ZZ}S_N - M_{ZN}S_Z
= \sum_{i,j} w_iw_j\bigl[a_i^2(1-a_j) - a_i(1-a_i)a_j\bigr]
= \sum_{i,j} w_iw_j\bigl[a_i^2 - a_ia_j\bigr].
\label{eq:dyedup2}
\end{equation}
%
Symmetrise in $(i, j)$:
%
\begin{equation}
= \tfrac12\sum_{i,j} w_iw_j\,(a_i-a_j)^2
\;=\; W^2\Var_w\!\left(\frac{Z}{A}\right),
\qquad W=\sum_i w_i,
\label{eq:dyedup3}
\end{equation}
%
with $\Var_w$ the $w$-normalised variance. Therefore
%
\begin{equation}
\boxed{\ \frac{\dd\Ye}{\dd\up}\bigg|_{\Sigma X=1}
\;=\; \frac{W^2}{kT\,S_N}\Var_w\!\left(\frac{Z}{A}\right) \;>\; 0\ }
\label{eq:monotonicity}
\end{equation}
\derivedhere

strictly, unless every abundant species has the same $Z/A$. $\square$

This is a satisfying result for three reasons. It \textbf{licenses the outer
bisection} (\code{if F[1] > 0: hi\_p = u\_p}) as a theorem rather than a hope.
It \textbf{explains the damping cap} (\S\ref{sec:newton}): the derivative is
proportional to the compositional spread in $Z/A$, so near a single-species
composition the map is nearly flat and Newton overshoots. And it
\textbf{identifies the exact degenerate case} --- a composition of species all
with identical $Z/A$ (e.g. pure $\alpha$-chain matter, $Z/A = 1/2$ throughout)
makes \Ye{} independent of \up{} entirely.

Be precise about what that degeneracy is, because the natural way to say it is
backwards. With every $Z/A = \tfrac12$, $\Ye \equiv 0.5$ \emph{identically},
whatever \up{} does. So $\Ye = 0.5$ is not the unachievable value --- it is the
\textbf{only} achievable one, and $F_1 \equiv 0$ leaves \up{} completely
undetermined (the Jacobian's second row vanishes); any target $\Ye \ne 0.5$ is
infeasible, and the bisection has no sign change to bracket. Both failure modes
--- an undetermined \up{} and an infeasible constraint --- are the same
vanishing $\Var_w(Z/A)$. That is not hypothetical: an $\alpha$-chain-only
network at exactly $\Ye = 0.5$ is the textbook NSE example, and it is precisely
where the solve is degenerate. The production networks avoid it by containing
odd-$A$ species.

\begin{warn}
\textbf{Honest caveat.} Claim 2 is proved \emph{along the $\sum X = 1$ curve for
the exact solution}. The fallback's outer loop calls an inner solve that
terminates on a finite tolerance, so the composed map is monotone only up to
that tolerance --- which is why the fallback still returns \code{converged=False}
with a residual rather than assuming success. The code is right not to trust the
proof further than it goes.
\end{warn}

% =====================================================================
\subsection{Batching, and why the scalar path survives}
\label{sec:batching}

\code{solve\_nse\_batch}\,/\,\code{solve\_qse\_batch} run one vectorised
damped-Newton loop over all states with per-row convergence masking and per-row
damping, then hand any row that did not converge in the Newton phase to the
\textbf{scalar} solver, which retries Newton and applies the bisection fallback.

Two design points worth stealing:

\begin{itemize}
\item \textbf{\code{\_batched\_2x2\_step}\,/\,\code{\_batched\_3x3\_step} use
  closed-form Cramer\,/\,adjugate solves, not \code{np.linalg.solve}.} The reason
  is in the docstring: a batched LAPACK solve \emph{aborts the whole batch} if
  any single row is singular. A closed-form solve never raises; singular rows
  are flagged and dropped to the fallback. Vectorisation must not make
  robustness worse than the loop it replaces.
\item \textbf{The batched result is required to be identical to a per-state
  loop}, and that is a test (\code{test\_nse\_batch\_matches\_scalar\_loop},
  \code{test\_qse\_batch\_matches\_scalar\_loop}), not a hope. Same
  \code{EXP\_CLIP}, same \code{\_MAX\_STEP\_MEV}, same \code{\_MAX\_ITER}, same
  tol, same float64.
\end{itemize}

% =====================================================================
\subsection{QSE: the cluster and \texorpdfstring{\uG}{uG}}
\label{sec:qsecluster}

Now the generalisation of \S\ref{sec:regimes}. Group $G$ (the silicon group) is
internally equilibrated but not equilibrated with the free nucleons; non-members
are in NSE with the nucleons. The chemical-potential ansatz gains one offset:
%
\begin{equation}
\mu_i \;=\; Z_i\mu_p + N_i\mu_n + \uG\,\mathbf{1}[i\in G],
\label{eq:qseansatz}
\end{equation}
%
so in the Saha exponent,

\begin{srclisting}
expo = logC + (Z*u_p + N*u_n)/kT
if group is not None:
    expo = expo + np.where(group, u_g/kT, 0.0)
\end{srclisting}

Three unknowns $(\up, \un, \uG)$ need three constraints; the third is the
\textbf{measured group mass fraction}:
%
\begin{equation}
F_2 \;=\; \sum_{i\in G} X_i - X_G^{\mathrm{meas}} \;=\; 0 .
\label{eq:qseF2}
\end{equation}
%
The $3\times3$ Jacobian is the same construction with
$\dd g = \mathbf{1}_G$:
%
\begin{equation}
J = \frac{1}{kT}\begin{pmatrix}
\sum XZ & \sum XN & \sum X\mathbf{1}_G\\
\sum XZ^2/A & \sum XZN/A & \sum XZ\mathbf{1}_G/A\\
\sum XZ\mathbf{1}_G & \sum XN\mathbf{1}_G & \sum X\mathbf{1}_G
\end{pmatrix}
\label{eq:qsejacobian}
\end{equation}
%
Note $J[0,2] = J[2,2]$: both are $\sum_{i\in G} X_i$, because
$\partial(\sum X)/\partial\uG$ and $\partial(X_G)/\partial\uG$ are the same sum.
A useful hand-check when reading the code.

\textbf{$\uG = 0$ recovers NSE}, and that is a test:
\code{test\_degenerates\_to\_nse\_when\_group\_mass\_is\_nse} --- feed the QSE
solver the group mass \emph{taken from the NSE solution}, and require
$|\uG| < 10^{-6}$\,MeV. \resultsref{2026-07-10} records it. A solver that could
not do this would be solving a different problem.

\begin{result}
\textbf{What \uG{} means physically.} $\exp(\uG/kT)$ is the factor by which the
whole silicon group is over- or under-populated relative to NSE at the same
nucleon potentials. Measured on pre-stall trajectory rows: median
\textbf{$+4.3$\,/\,$+4.1$\,MeV} \resultsref{2026-07-10}. At
$kT \approx 0.35$\,MeV that is $\exp(12) \approx 2\times10^{5}$ --- the Si group
is over five orders of magnitude \emph{over}-populated relative to NSE, which is
exactly what ``mid-burn, not yet at the iron peak'' means quantitatively.
Reading \uG{} is the cleanest single number for ``how far through silicon
burning is this state''.
\end{result}

% =====================================================================
\subsection{QSE as a slow manifold}
\label{sec:slowmanifold}

The reason QSE is a \emph{dynamical} statement and not just an algebraic one.

Split the RHS by timescale:
%
\begin{equation}
\dot Y \;=\; \frac{1}{\epsilon}f_{\mathrm{fast}}(Y) + f_{\mathrm{slow}}(Y),
\qquad \epsilon \ll 1,
\label{eq:timescalesplit}
\end{equation}
%
where $f_{\rm fast}$ collects the near-balanced strong pairs and $f_{\rm slow}$
the bridges and weak columns. Singular perturbation theory
\citep[geometric form][]{tikhonov1952, fenichel1979} says: provided the fast
subsystem's root is asymptotically stable (Tikhonov's hypothesis --- physically
satisfied for relaxation toward detailed balance), from generic initial
conditions the system relaxes on timescale $\epsilon$ onto the manifold
$\Mslow = \{f_{\rm fast}(Y) = 0\}$, and thereafter moves \emph{within} \Mslow{}
at the rate set by $f_{\rm slow}$. The same structure under the names ILDM and
CSP is the workhorse of combustion kinetics reduction \citep{maas1992, lam1994}.

\begin{itemize}
\item \textbf{\Mslow{} is exactly the QSE manifold.} $f_{\rm fast} = 0$ is
  ``every intra-cluster reaction balances'', which is the chemical-potential
  ansatz of \S\ref{sec:qsecluster}.
\item \textbf{Its dimension is the QSE parameter count}, $2 + g$ (plus \Ye),
  which is why the reduction works at all.
\item \textbf{The stiffness ratio is $1/\epsilon$.} So \emph{stiffness and QSE
  are the same phenomenon} --- Tier~0 \S V.3 says this, and here is the reason:
  a stiff system is one with an attracting low-dimensional manifold, and the
  manifold is the equilibrium.
\item \textbf{The group mass is the slow coordinate.} This is why QSE needs a
  measured $X_G$: the reduction removes the fast directions and leaves the slow
  ones to be supplied dynamically.
\end{itemize}

Two consequences that carry into \S\ref{part:integration} and Tier~4:

\begin{enumerate}
\item An implicit solver's step size is limited by the \emph{slow} dynamics once
  on the manifold; an explicit solver's by the \emph{fast} ones forever. That is
  the whole argument for BDF (\S\ref{sec:bdf}).
\item An emulator learning one $\Delta t$ step is learning the \textbf{flow map
  restricted to \Mslow} for states already on it, and the \emph{relaxation onto}
  \Mslow{} for states off it. Those are different functions with different
  smoothness. The training grid consists overwhelmingly of the second kind
  (\S\ref{sec:vacuous} --- forward, and the last forward dependency this section
  has: in the source that measurement precedes this one); deployment is the
  first. Naming this makes the distribution-shift problem concrete rather than
  vague.
\end{enumerate}

% =====================================================================
\subsection{The group boundary is a config variant, not a fact}
\label{sec:groupboundary}

\code{configs/qse\_groups.yaml}:

\begin{srclisting}
default: {a_min: 24, a_max: 45, exclude_light: true}   # Hix & Thielemann span
a28_up:  {a_min: 28,            exclude_light: true}
a24_46:  {a_min: 24, a_max: 46, exclude_light: true}   # Sc45 inside the group
\end{srclisting}

\code{exclude\_light: true} forces $Z > 2$ --- free $n$, $p$, \nuc{He}{4} are
never group members (they are the reservoir the group exchanges \emph{with};
putting them inside would make the cluster the whole network).
\code{load\_group\_mask} additionally refuses any variant that excludes
\nuc{Si}{28}, the canonical anchor.

Membership counts: 29\,/\,50 (\code{default}, \msn{80}\,/\,\msn{151}),
40\,/\,105 (\code{a28\_up}). \resultsref{2026-07-10}

\paragraph{Why the boundary matters and is not resolvable by argument.} A
reaction is a \emph{bridge} iff it crosses the boundary.
\nuc{Sc}{45}$(p,\gamma)$\nuc{Ti}{46} --- the literature high-\Ye{} bottleneck
\citep[\S VIb]{wac1973}; \citep{ht1996} --- is a boundary crossing under
\code{a24\_46} and an internal reaction under \code{default}. Under
\code{a24\_46} it ranks \textbf{2/251} inter-group carriers (share 0.10--0.12)
on relaxed high-\Ye{} QSE-window rows; under \code{default} its \emph{feeder}
\code{Ca44(p,g)Sc45} ranks \#4 instead. \resultsref{2026-07-10\,/\,2026-07-12}

\code{CLAUDE.md}'s rule --- every group\,/\,bridge\,/\,inter-group analysis
reports \textbf{both} variants, and a verdict that flips between them is a
\emph{measured decision} --- is the correct response to an unfalsifiable
modelling choice. Measured outcome: the boundary variant changes top-$k$
concentration by $\le 0.02$, so no verdict depends on it
\resultsref{2026-07-12}. That is the result you want from a sensitivity variant:
it was worth checking and it did not matter.

% =====================================================================
\subsection{What was measured: the solver is right}
\label{sec:solververify}

S9's headline outputs are negative, and each is informative in a different way.

\begin{result}
\textbf{(1) The solver is right.} NSE cross-check vs pynucastro on the canonical
27-state grid: max $|\Delta\log_{10}X|$ over $X > 10^{-10}$ is
\textbf{$2.7\times10^{-10}$} (\msn{80})\,/\,\textbf{$1.5\times10^{-10}$}
(\msn{151}) against a $10^{-6}$ gate; 27/27 states, all Newton, no fallback
needed. \resultsref{2026-07-10} Two independent solvers on shared data agreeing
to ten digits is as strong as this kind of check gets.
\end{result}

\begin{aside}
\textbf{Where the other two went, and why.} Results \textbf{(2)} the
single-Si-cluster plateau is not confirmed, and \textbf{(3)} the maskable set is
empty are both stated against \kap{} and against the departure diagnostics
\kap{} is compared with, so this report derives them in
\S\ref{sec:negresults}, after the flux algebra that gives them meaning. The
source states all three together at \S III.12; splitting them is this report's
editorial choice, not the source's, and the reading it suggests --- that
everything in \emph{this} section is what makes the other two legible, since
without a solver verified to ten digits ``the labels' equilibria are displaced''
would be hard to tell from ``our reference is wrong'' --- is likewise the
report's, not a claim the source makes.
\end{aside}

% =====================================================================
\subsection{Code and oracles}
\label{sec:codes9}

\textbf{Read:} \code{qse/coeffs.py} (132 lines --- the derivation of
\S\ref{sec:saha} as code), \code{qse/solver.py} (581 lines: scalar NSE $\to$
scalar QSE $\to$ batched, in that order), \code{qse/diagnostics.py} (85 lines
--- small and load-bearing).

\textbf{Oracles:}

\begin{center}
\begin{tabularx}{\textwidth}{@{}L{7.0cm} Y@{}}
\toprule
\textbf{test} & \textbf{pins} \\
\midrule
\code{test\_qse\_coeffs.py::test\_mass\_fractions\_match\_at\_fixed\_u} &
  \S\ref{sec:saha}, $\le 10^{-12}$ rel, solver excluded \\
\code{test\_qse\_solver.py::TestNSECrossCheck::test\_matches\_pyna\_solver} &
  the converged solution, 3 box states \\
\code{test\_constraints\_satisfied} & $F_0$, $F_1$ at the returned solution \\
\code{test\_degenerates\_to\_nse\_when\_group\_mass\_is\_nse} &
  $\uG \to 0$ (\S\ref{sec:qsecluster}) \\
\code{test\_offset\_group\_mass\_moves\_u\_group} &
  the converse --- \uG{} responds \\
\code{test\_group\_mask\_contains\_si28\_not\_light} &
  the config contract (\S\ref{sec:groupboundary}) \\
\code{test\_nse\_batch\_matches\_scalar\_loop},
\code{test\_qse\_batch\_matches\_scalar\_loop} & \S\ref{sec:batching} \\
\code{test\_weak\_columns\_structurally\_excluded} &
  \textbf{invariant \#2} (\S\ref{sec:departure}) \\
\code{test\_reaction\_delta\_max\_over\_participants} &
  the conservative $\delta_r$ (\S\ref{sec:departure}) \\
\bottomrule
\end{tabularx}
\end{center}

\begin{repopointer}
\textbf{\textcircled{4} SEE:} notebook \code{06-qse-bridges.ipynb} --- Fig 1 the
Si group under both boundary conventions, Fig 2 inter-group carriers side by
side, Fig 3 whether the boundary changes the answer (it does not).
\end{repopointer}

% =====================================================================
\subsection{Self-check}
\label{sec:selfcheck-s9}

\begin{selfcheck}
\item Show that nuclei are non-degenerate at box centre while electrons are not.
  Which quantity differs, and by how much?
\item Derive $u_i = Z_i \up + N_i \un + B_i$ and say where the binding energy
  came from.
\item Write $\log C_i$ and identify each of its six terms. Why is $A_{\rm nuc}$
  used in one place and integer $A$ in another?
\item Prove $\dd\Ye/\dd\up > 0$ along $\sum X = 1$. What does the result say
  about when the solve is degenerate, and about why the Newton step is capped?
\item QSE has three unknowns. Name them and name the three constraints. Which
  one is dynamical rather than thermodynamic?
\item Explain ``stiffness and QSE are the same phenomenon'' in terms of a slow
  manifold.
\end{selfcheck}

\begin{aside}
The source's remaining two S9 self-check questions --- $\delta$ vs \rQSE{}, and
the empty maskable set against \kap's continuum --- test material that this
report reaches in \S\ref{part:fluxes}, so they are asked there
(\S\ref{sec:selfcheck-s8}). Appendix~\ref{app:concordance} records the move.
\end{aside}
